% =====================================================================
% Seção de Método -- MEST (condição L2d) -- VERSÃO ESTENDIDA (~2,7 páginas SBC)
% Padrão SBC: % =====================================================================
% Seção de Método -- MEST (condição L2d)
% Padrão SBC: % =====================================================================
% Seção de Método -- MEST (condição L2d)
% Padrão SBC: % =====================================================================
% Seção de Método -- MEST (condição L2d)
% Padrão SBC: \input{secao_metodo_MEST} no corpo do artigo.
% Pacotes assumidos: amsmath, amssymb.
% =====================================================================

\section{MEST: Memória Episódica de Slots Tipados}
\label{sec:metodo}

A ideia do método cabe em uma frase: \emph{uma pergunta sobre uma conversa passada
já chega estruturada, e a memória deve ser guardada com a mesma estrutura}. Quando
alguém pergunta ``o que o James adotou em abril de 2022?'', a pergunta nomeia uma
pessoa (James), uma ação (adotar) e uma data (abril de 2022), e deixa exatamente um
campo em branco --- o objeto ---, que é justamente o que ela quer saber. Perguntas
de memória conversacional são, quase sempre, uma tupla \textbf{(quem, fez-o-quê,
com-o-quê, quando)} com um dos campos vazio. Um índice comum de palavras não
enxerga essa estrutura: ele só sabe responder ``quais trechos mencionam a palavra
$X$''. O \textbf{MEST} (\emph{Memória Episódica de Slots Tipados}) guarda cada
trecho da conversa já anotado com \emph{quem} falou, \emph{que ação} ocorreu,
\emph{sobre o quê}, de \emph{que tipo} de coisa se trata e \emph{quando} --- e
responde a pergunta procurando quem completa o campo que ficou vazio.

Em resumo, o método faz três movimentos. \textbf{(1)~Recorta} cada conversa em
\emph{episódios}: pequenos blocos de turnos consecutivos que tratam do mesmo
assunto, sendo o menor deles o par ``pergunta e resposta'' do diálogo.
\textbf{(2)~Pendura} cada episódio em cinco tipos de rótulo, os
\emph{slots}: ator, predicado, conceito, tipo e tempo. \textbf{(3)~Traduz} a pergunta para
esses mesmos cinco rótulos e pontua os episódios pelos rótulos que eles têm em
comum com a pergunta, somando ainda um canal de similaridade textual. Duas
propriedades atravessam os três movimentos: a memória é construída \emph{sem
nenhuma chamada a modelo de linguagem}, e \emph{nenhum limiar foi escolhido
olhando as respostas certas} --- todos são estimados do próprio corpus
(Seção~\ref{sec:metodo:calibracao}).

\subsection{Como a memória é construída}
\label{sec:metodo:substrato}

Uma única passagem de análise sintática por turno entrega de uma só vez quatro dos
cinco slots: sintagmas nominais (\textsf{conceito}), lemas de verbos de conteúdo
(\textsf{predicado}), menções a pessoas (\textsf{ator}) e datas (\textsf{tempo},
resolvidas contra a data da sessão e indexadas por ano, mês e dia). O quinto vem de elevar cada conceito aos seus hiperônimos na WordNet: é o
slot \textsf{tipo}, que faz ``quais \emph{comidas}'' alcançar um turno que diz
``\emph{frango assado}''.

O resultado é um grafo bipartido: de um lado os episódios, do outro os slots, e uma
aresta sempre que aquele rótulo foi \emph{literalmente pronunciado} naquele
episódio --- nada é inventado nem reescrito. Sobre esse grafo colocamos uma \emph{rotação}
(estrutura de \emph{fatgraph}): uma ordem cíclica fixa das arestas em torno de cada
vértice, $G=(V,H,\alpha,\sigma)$. A rotação é projetada, e rende dois objetos úteis. Em torno de um \textbf{slot}, a ordem é cronológica: percorrer
essa órbita é ler ``tudo o que esta pessoa, este verbo ou este tópico tocou, em
ordem'', sem precisar ranquear nada. Em torno de um \textbf{episódio}, a ordem é
sempre $\textsf{ator}\prec\textsf{predicado}\prec\textsf{conceito}\prec
\textsf{tipo}\prec\textsf{tempo}$, e portanto os pares vizinhos --- os
\emph{cantos} --- são exatamente (quem, fez-o-quê), (fez-o-quê, com-o-quê) e
(com-o-quê, quando). É a mesma tupla do primeiro parágrafo: uma pergunta é um canto
com um lado em branco.

\subsection{Como uma pergunta é respondida}
\label{sec:metodo:consulta}

A pergunta passa pelo mesmo extrator, o que produz o conjunto $\Lambda(q)$ dos
slots que ela nomeia e que existem na memória. Cada slot encontrado ``acende'' os
episódios ligados a ele, mas não com o mesmo peso: um rótulo que aparece em $200$
episódios não distingue nenhum deles, ao passo que um rótulo presente em três
praticamente resolve a busca sozinho. Por isso o peso de um slot cai com o seu
grau, e acima de um corte $h_\kappa$ ele deixa de apontar episódios e vira só um
bônus fixo --- a regra é ``\emph{um rótulo muito comum serve de filtro, nunca de
ponte}''. Somando os canais estruturais com um canal de similaridade textual, o
escore de um episódio $\varepsilon$ é
\begin{equation}
\begin{aligned}
\psi(v)&=w_{\kappa(v)}\bigl(1+\ln \deg(v)\bigr)^{-\gamma}\ \ \text{se}\ \deg(v)<h_{\kappa(v)};\quad
\psi(v)=w_{H}\ \text{caso contrário},\\
s(\varepsilon)&=w_{D}\,\delta(q,\varepsilon)+\sum_{v\in\Lambda^{*}(q,\varepsilon)}\psi(v)
+\beta\,\mathbb{1}\bigl[\varepsilon\in\mathrm{Orb}(v^{\star})\bigr],
\end{aligned}
\label{eq:psi}
\end{equation}
em que $\Lambda^{*}(q,\varepsilon)$ são os slots da pergunta (exceto os de ator)
ligados a $\varepsilon$, e $\delta(q,\varepsilon)$ é o maior cosseno entre a
pergunta e o episódio ou qualquer um de seus turnos.

O termo $\beta$ trata das perguntas que pedem uma \emph{lista} (``quais comidas
ela gosta?''). Aqui a estrutura de rotação faz algo que um buscador denso não faz:
como a órbita de um slot já \emph{é} a lista cronológica de tudo o que ele tocou,
basta tomar o slot mais raro entre os que a pergunta nomeou e promover a órbita
inteira --- sem segunda busca nem reordenação.

O ator recebe tratamento diferente dos demais, e a razão é simples: como quase toda
pergunta nomeia um dos falantes, somar um bônus por ator empurraria metade do grafo
para cima e não separaria nada. O ator, portanto, não é um canal, é um
\emph{filtro}: $\hat{s}(\varepsilon)=s(\varepsilon)\,[\pi+(1-\pi)
\min(\omega(a,\varepsilon)/\Omega,1)]$, onde $\omega(a,\varepsilon)$ mede quanto do
conteúdo daquele episódio saiu da boca de $a$. Um episódio é um par de turnos e
contém, por construção, fala dos dois participantes; então ``$a$ falou aqui'' não
diz nada, mas ``$a$ disse a maior parte disto'' diz.

Uma última assimetria é proposital: quem torna um trecho \emph{encontrável} é o
episódio, mas quem ocupa o orçamento de \emph{tokens} é o turno. Entregar episódios inteiros gastaria o
orçamento com poucos blocos grandes e perderia amplitude; entregar turnos soltos,
por outro lado, quebraria justamente o par pergunta--resposta que motivou o
episódio. A solução é pontuar turnos, deixando que cada um herde o que o seu
episódio conquistou:
\begin{equation}
s(t)=\hat{s}\bigl(\varepsilon(t)\bigr)+w_{D}\cos(q,x_t)
+\lambda\,w_{D}\!\!\max_{t'\in\varepsilon(t),\,t'\neq t}\!\!\cos(q,x_{t'}).
\label{eq:turno}
\end{equation}
O último termo é a regra do par escrita como aritmética: um turno herda uma fração
$\lambda$ da melhor semelhança entre os seus vizinhos de episódio. Assim, a
resposta ``era uma peça contemporânea chamada \emph{Finding Freedom}'' sobe junto
com o turno que menciona o assunto, mesmo não tendo palavra alguma em comum com a
pergunta. Tomam-se então os turnos em ordem decrescente até esgotar o orçamento.

\subsection{De onde vêm os limiares}
\label{sec:metodo:calibracao}

Resta uma pergunta incômoda: de onde saem $h_\kappa$, $\pi$ e os demais números?
Ajustá-los contra o conjunto anotado é prática comum, e é também o que impede levar
um método a outro corpus. Adotamos um critério simples --- \emph{o parâmetro
precisa das respostas certas para ser fixado?} --- e trocamos cada constante por um
estimador medido, durante a construção, sobre o corpus \emph{não} anotado. O corte de rótulo comum vira o quantil $0{,}99$ dos graus
daquele tipo, $h_\kappa=\lceil Q_{0,99}\{\deg(v):v\in\mathcal{S}_\kappa\}\rceil$:
uma contagem absoluta quebraria com a escala, pois num corpus dez vezes maior todo
slot a ultrapassaria e o mecanismo se desligaria sozinho. Pelo mesmo raciocínio, o
limiar de paráfrase vira um quantil alto dos cossenos realmente observados --- é
propriedade do codificador, não da tarefa --- e o filtro de ator vira
$\pi=1/|\Sigma|$ e $\Omega=\mathrm{mediana}_\varepsilon\max_a\omega(a,\varepsilon)$,
apertando-se sozinho quando há mais falantes. Já as palavras que só servem de
moldura (``conversa'', ``data'', ``tipo'') passam a ser detectadas por um contraste
de frequências, o análogo do IDF do lado da pergunta: uma palavra de tópico é
frequente nas perguntas \emph{porque} o é nas conversas, ao passo que uma de
moldura é frequente nas perguntas e ausente do que as pessoas disseram. Os pesos $w_\kappa$, $\lambda$ e $\gamma$ seguem fixos de propósito: afirmam algo
sobre o \emph{modelo} --- que um casamento de conceito informa mais do que um
palpite de hiperônimo ---, não sobre o corpus.
 no corpo do artigo.
% Pacotes assumidos: amsmath, amssymb.
% =====================================================================

\section{MEST: Memória Episódica de Slots Tipados}
\label{sec:metodo}

A ideia do método cabe em uma frase: \emph{uma pergunta sobre uma conversa passada
já chega estruturada, e a memória deve ser guardada com a mesma estrutura}. Quando
alguém pergunta ``o que o James adotou em abril de 2022?'', a pergunta nomeia uma
pessoa (James), uma ação (adotar) e uma data (abril de 2022), e deixa exatamente um
campo em branco --- o objeto ---, que é justamente o que ela quer saber. Perguntas
de memória conversacional são, quase sempre, uma tupla \textbf{(quem, fez-o-quê,
com-o-quê, quando)} com um dos campos vazio. Um índice comum de palavras não
enxerga essa estrutura: ele só sabe responder ``quais trechos mencionam a palavra
$X$''. O \textbf{MEST} (\emph{Memória Episódica de Slots Tipados}) guarda cada
trecho da conversa já anotado com \emph{quem} falou, \emph{que ação} ocorreu,
\emph{sobre o quê}, de \emph{que tipo} de coisa se trata e \emph{quando} --- e
responde a pergunta procurando quem completa o campo que ficou vazio.

Em resumo, o método faz três movimentos. \textbf{(1)~Recorta} cada conversa em
\emph{episódios}: pequenos blocos de turnos consecutivos que tratam do mesmo
assunto, sendo o menor deles o par ``pergunta e resposta'' do diálogo.
\textbf{(2)~Pendura} cada episódio em cinco tipos de rótulo, os
\emph{slots}: ator, predicado, conceito, tipo e tempo. \textbf{(3)~Traduz} a pergunta para
esses mesmos cinco rótulos e pontua os episódios pelos rótulos que eles têm em
comum com a pergunta, somando ainda um canal de similaridade textual. Duas
propriedades atravessam os três movimentos: a memória é construída \emph{sem
nenhuma chamada a modelo de linguagem}, e \emph{nenhum limiar foi escolhido
olhando as respostas certas} --- todos são estimados do próprio corpus
(Seção~\ref{sec:metodo:calibracao}).

\subsection{Como a memória é construída}
\label{sec:metodo:substrato}

Uma única passagem de análise sintática por turno entrega de uma só vez quatro dos
cinco slots: sintagmas nominais (\textsf{conceito}), lemas de verbos de conteúdo
(\textsf{predicado}), menções a pessoas (\textsf{ator}) e datas (\textsf{tempo},
resolvidas contra a data da sessão e indexadas por ano, mês e dia). O quinto vem de elevar cada conceito aos seus hiperônimos na WordNet: é o
slot \textsf{tipo}, que faz ``quais \emph{comidas}'' alcançar um turno que diz
``\emph{frango assado}''.

O resultado é um grafo bipartido: de um lado os episódios, do outro os slots, e uma
aresta sempre que aquele rótulo foi \emph{literalmente pronunciado} naquele
episódio --- nada é inventado nem reescrito. Sobre esse grafo colocamos uma \emph{rotação}
(estrutura de \emph{fatgraph}): uma ordem cíclica fixa das arestas em torno de cada
vértice, $G=(V,H,\alpha,\sigma)$. A rotação é projetada, e rende dois objetos úteis. Em torno de um \textbf{slot}, a ordem é cronológica: percorrer
essa órbita é ler ``tudo o que esta pessoa, este verbo ou este tópico tocou, em
ordem'', sem precisar ranquear nada. Em torno de um \textbf{episódio}, a ordem é
sempre $\textsf{ator}\prec\textsf{predicado}\prec\textsf{conceito}\prec
\textsf{tipo}\prec\textsf{tempo}$, e portanto os pares vizinhos --- os
\emph{cantos} --- são exatamente (quem, fez-o-quê), (fez-o-quê, com-o-quê) e
(com-o-quê, quando). É a mesma tupla do primeiro parágrafo: uma pergunta é um canto
com um lado em branco.

\subsection{Como uma pergunta é respondida}
\label{sec:metodo:consulta}

A pergunta passa pelo mesmo extrator, o que produz o conjunto $\Lambda(q)$ dos
slots que ela nomeia e que existem na memória. Cada slot encontrado ``acende'' os
episódios ligados a ele, mas não com o mesmo peso: um rótulo que aparece em $200$
episódios não distingue nenhum deles, ao passo que um rótulo presente em três
praticamente resolve a busca sozinho. Por isso o peso de um slot cai com o seu
grau, e acima de um corte $h_\kappa$ ele deixa de apontar episódios e vira só um
bônus fixo --- a regra é ``\emph{um rótulo muito comum serve de filtro, nunca de
ponte}''. Somando os canais estruturais com um canal de similaridade textual, o
escore de um episódio $\varepsilon$ é
\begin{equation}
\begin{aligned}
\psi(v)&=w_{\kappa(v)}\bigl(1+\ln \deg(v)\bigr)^{-\gamma}\ \ \text{se}\ \deg(v)<h_{\kappa(v)};\quad
\psi(v)=w_{H}\ \text{caso contrário},\\
s(\varepsilon)&=w_{D}\,\delta(q,\varepsilon)+\sum_{v\in\Lambda^{*}(q,\varepsilon)}\psi(v)
+\beta\,\mathbb{1}\bigl[\varepsilon\in\mathrm{Orb}(v^{\star})\bigr],
\end{aligned}
\label{eq:psi}
\end{equation}
em que $\Lambda^{*}(q,\varepsilon)$ são os slots da pergunta (exceto os de ator)
ligados a $\varepsilon$, e $\delta(q,\varepsilon)$ é o maior cosseno entre a
pergunta e o episódio ou qualquer um de seus turnos.

O termo $\beta$ trata das perguntas que pedem uma \emph{lista} (``quais comidas
ela gosta?''). Aqui a estrutura de rotação faz algo que um buscador denso não faz:
como a órbita de um slot já \emph{é} a lista cronológica de tudo o que ele tocou,
basta tomar o slot mais raro entre os que a pergunta nomeou e promover a órbita
inteira --- sem segunda busca nem reordenação.

O ator recebe tratamento diferente dos demais, e a razão é simples: como quase toda
pergunta nomeia um dos falantes, somar um bônus por ator empurraria metade do grafo
para cima e não separaria nada. O ator, portanto, não é um canal, é um
\emph{filtro}: $\hat{s}(\varepsilon)=s(\varepsilon)\,[\pi+(1-\pi)
\min(\omega(a,\varepsilon)/\Omega,1)]$, onde $\omega(a,\varepsilon)$ mede quanto do
conteúdo daquele episódio saiu da boca de $a$. Um episódio é um par de turnos e
contém, por construção, fala dos dois participantes; então ``$a$ falou aqui'' não
diz nada, mas ``$a$ disse a maior parte disto'' diz.

Uma última assimetria é proposital: quem torna um trecho \emph{encontrável} é o
episódio, mas quem ocupa o orçamento de \emph{tokens} é o turno. Entregar episódios inteiros gastaria o
orçamento com poucos blocos grandes e perderia amplitude; entregar turnos soltos,
por outro lado, quebraria justamente o par pergunta--resposta que motivou o
episódio. A solução é pontuar turnos, deixando que cada um herde o que o seu
episódio conquistou:
\begin{equation}
s(t)=\hat{s}\bigl(\varepsilon(t)\bigr)+w_{D}\cos(q,x_t)
+\lambda\,w_{D}\!\!\max_{t'\in\varepsilon(t),\,t'\neq t}\!\!\cos(q,x_{t'}).
\label{eq:turno}
\end{equation}
O último termo é a regra do par escrita como aritmética: um turno herda uma fração
$\lambda$ da melhor semelhança entre os seus vizinhos de episódio. Assim, a
resposta ``era uma peça contemporânea chamada \emph{Finding Freedom}'' sobe junto
com o turno que menciona o assunto, mesmo não tendo palavra alguma em comum com a
pergunta. Tomam-se então os turnos em ordem decrescente até esgotar o orçamento.

\subsection{De onde vêm os limiares}
\label{sec:metodo:calibracao}

Resta uma pergunta incômoda: de onde saem $h_\kappa$, $\pi$ e os demais números?
Ajustá-los contra o conjunto anotado é prática comum, e é também o que impede levar
um método a outro corpus. Adotamos um critério simples --- \emph{o parâmetro
precisa das respostas certas para ser fixado?} --- e trocamos cada constante por um
estimador medido, durante a construção, sobre o corpus \emph{não} anotado. O corte de rótulo comum vira o quantil $0{,}99$ dos graus
daquele tipo, $h_\kappa=\lceil Q_{0,99}\{\deg(v):v\in\mathcal{S}_\kappa\}\rceil$:
uma contagem absoluta quebraria com a escala, pois num corpus dez vezes maior todo
slot a ultrapassaria e o mecanismo se desligaria sozinho. Pelo mesmo raciocínio, o
limiar de paráfrase vira um quantil alto dos cossenos realmente observados --- é
propriedade do codificador, não da tarefa --- e o filtro de ator vira
$\pi=1/|\Sigma|$ e $\Omega=\mathrm{mediana}_\varepsilon\max_a\omega(a,\varepsilon)$,
apertando-se sozinho quando há mais falantes. Já as palavras que só servem de
moldura (``conversa'', ``data'', ``tipo'') passam a ser detectadas por um contraste
de frequências, o análogo do IDF do lado da pergunta: uma palavra de tópico é
frequente nas perguntas \emph{porque} o é nas conversas, ao passo que uma de
moldura é frequente nas perguntas e ausente do que as pessoas disseram. Os pesos $w_\kappa$, $\lambda$ e $\gamma$ seguem fixos de propósito: afirmam algo
sobre o \emph{modelo} --- que um casamento de conceito informa mais do que um
palpite de hiperônimo ---, não sobre o corpus.
 no corpo do artigo.
% Pacotes assumidos: amsmath, amssymb.
% =====================================================================

\section{MEST: Memória Episódica de Slots Tipados}
\label{sec:metodo}

A ideia do método cabe em uma frase: \emph{uma pergunta sobre uma conversa passada
já chega estruturada, e a memória deve ser guardada com a mesma estrutura}. Quando
alguém pergunta ``o que o James adotou em abril de 2022?'', a pergunta nomeia uma
pessoa (James), uma ação (adotar) e uma data (abril de 2022), e deixa exatamente um
campo em branco --- o objeto ---, que é justamente o que ela quer saber. Perguntas
de memória conversacional são, quase sempre, uma tupla \textbf{(quem, fez-o-quê,
com-o-quê, quando)} com um dos campos vazio. Um índice comum de palavras não
enxerga essa estrutura: ele só sabe responder ``quais trechos mencionam a palavra
$X$''. O \textbf{MEST} (\emph{Memória Episódica de Slots Tipados}) guarda cada
trecho da conversa já anotado com \emph{quem} falou, \emph{que ação} ocorreu,
\emph{sobre o quê}, de \emph{que tipo} de coisa se trata e \emph{quando} --- e
responde a pergunta procurando quem completa o campo que ficou vazio.

Em resumo, o método faz três movimentos. \textbf{(1)~Recorta} cada conversa em
\emph{episódios}: pequenos blocos de turnos consecutivos que tratam do mesmo
assunto, sendo o menor deles o par ``pergunta e resposta'' do diálogo.
\textbf{(2)~Pendura} cada episódio em cinco tipos de rótulo, os
\emph{slots}: ator, predicado, conceito, tipo e tempo. \textbf{(3)~Traduz} a pergunta para
esses mesmos cinco rótulos e pontua os episódios pelos rótulos que eles têm em
comum com a pergunta, somando ainda um canal de similaridade textual. Duas
propriedades atravessam os três movimentos: a memória é construída \emph{sem
nenhuma chamada a modelo de linguagem}, e \emph{nenhum limiar foi escolhido
olhando as respostas certas} --- todos são estimados do próprio corpus
(Seção~\ref{sec:metodo:calibracao}).

\subsection{Como a memória é construída}
\label{sec:metodo:substrato}

Uma única passagem de análise sintática por turno entrega de uma só vez quatro dos
cinco slots: sintagmas nominais (\textsf{conceito}), lemas de verbos de conteúdo
(\textsf{predicado}), menções a pessoas (\textsf{ator}) e datas (\textsf{tempo},
resolvidas contra a data da sessão e indexadas por ano, mês e dia). O quinto vem de elevar cada conceito aos seus hiperônimos na WordNet: é o
slot \textsf{tipo}, que faz ``quais \emph{comidas}'' alcançar um turno que diz
``\emph{frango assado}''.

O resultado é um grafo bipartido: de um lado os episódios, do outro os slots, e uma
aresta sempre que aquele rótulo foi \emph{literalmente pronunciado} naquele
episódio --- nada é inventado nem reescrito. Sobre esse grafo colocamos uma \emph{rotação}
(estrutura de \emph{fatgraph}): uma ordem cíclica fixa das arestas em torno de cada
vértice, $G=(V,H,\alpha,\sigma)$. A rotação é projetada, e rende dois objetos úteis. Em torno de um \textbf{slot}, a ordem é cronológica: percorrer
essa órbita é ler ``tudo o que esta pessoa, este verbo ou este tópico tocou, em
ordem'', sem precisar ranquear nada. Em torno de um \textbf{episódio}, a ordem é
sempre $\textsf{ator}\prec\textsf{predicado}\prec\textsf{conceito}\prec
\textsf{tipo}\prec\textsf{tempo}$, e portanto os pares vizinhos --- os
\emph{cantos} --- são exatamente (quem, fez-o-quê), (fez-o-quê, com-o-quê) e
(com-o-quê, quando). É a mesma tupla do primeiro parágrafo: uma pergunta é um canto
com um lado em branco.

\subsection{Como uma pergunta é respondida}
\label{sec:metodo:consulta}

A pergunta passa pelo mesmo extrator, o que produz o conjunto $\Lambda(q)$ dos
slots que ela nomeia e que existem na memória. Cada slot encontrado ``acende'' os
episódios ligados a ele, mas não com o mesmo peso: um rótulo que aparece em $200$
episódios não distingue nenhum deles, ao passo que um rótulo presente em três
praticamente resolve a busca sozinho. Por isso o peso de um slot cai com o seu
grau, e acima de um corte $h_\kappa$ ele deixa de apontar episódios e vira só um
bônus fixo --- a regra é ``\emph{um rótulo muito comum serve de filtro, nunca de
ponte}''. Somando os canais estruturais com um canal de similaridade textual, o
escore de um episódio $\varepsilon$ é
\begin{equation}
\begin{aligned}
\psi(v)&=w_{\kappa(v)}\bigl(1+\ln \deg(v)\bigr)^{-\gamma}\ \ \text{se}\ \deg(v)<h_{\kappa(v)};\quad
\psi(v)=w_{H}\ \text{caso contrário},\\
s(\varepsilon)&=w_{D}\,\delta(q,\varepsilon)+\sum_{v\in\Lambda^{*}(q,\varepsilon)}\psi(v)
+\beta\,\mathbb{1}\bigl[\varepsilon\in\mathrm{Orb}(v^{\star})\bigr],
\end{aligned}
\label{eq:psi}
\end{equation}
em que $\Lambda^{*}(q,\varepsilon)$ são os slots da pergunta (exceto os de ator)
ligados a $\varepsilon$, e $\delta(q,\varepsilon)$ é o maior cosseno entre a
pergunta e o episódio ou qualquer um de seus turnos.

O termo $\beta$ trata das perguntas que pedem uma \emph{lista} (``quais comidas
ela gosta?''). Aqui a estrutura de rotação faz algo que um buscador denso não faz:
como a órbita de um slot já \emph{é} a lista cronológica de tudo o que ele tocou,
basta tomar o slot mais raro entre os que a pergunta nomeou e promover a órbita
inteira --- sem segunda busca nem reordenação.

O ator recebe tratamento diferente dos demais, e a razão é simples: como quase toda
pergunta nomeia um dos falantes, somar um bônus por ator empurraria metade do grafo
para cima e não separaria nada. O ator, portanto, não é um canal, é um
\emph{filtro}: $\hat{s}(\varepsilon)=s(\varepsilon)\,[\pi+(1-\pi)
\min(\omega(a,\varepsilon)/\Omega,1)]$, onde $\omega(a,\varepsilon)$ mede quanto do
conteúdo daquele episódio saiu da boca de $a$. Um episódio é um par de turnos e
contém, por construção, fala dos dois participantes; então ``$a$ falou aqui'' não
diz nada, mas ``$a$ disse a maior parte disto'' diz.

Uma última assimetria é proposital: quem torna um trecho \emph{encontrável} é o
episódio, mas quem ocupa o orçamento de \emph{tokens} é o turno. Entregar episódios inteiros gastaria o
orçamento com poucos blocos grandes e perderia amplitude; entregar turnos soltos,
por outro lado, quebraria justamente o par pergunta--resposta que motivou o
episódio. A solução é pontuar turnos, deixando que cada um herde o que o seu
episódio conquistou:
\begin{equation}
s(t)=\hat{s}\bigl(\varepsilon(t)\bigr)+w_{D}\cos(q,x_t)
+\lambda\,w_{D}\!\!\max_{t'\in\varepsilon(t),\,t'\neq t}\!\!\cos(q,x_{t'}).
\label{eq:turno}
\end{equation}
O último termo é a regra do par escrita como aritmética: um turno herda uma fração
$\lambda$ da melhor semelhança entre os seus vizinhos de episódio. Assim, a
resposta ``era uma peça contemporânea chamada \emph{Finding Freedom}'' sobe junto
com o turno que menciona o assunto, mesmo não tendo palavra alguma em comum com a
pergunta. Tomam-se então os turnos em ordem decrescente até esgotar o orçamento.

\subsection{De onde vêm os limiares}
\label{sec:metodo:calibracao}

Resta uma pergunta incômoda: de onde saem $h_\kappa$, $\pi$ e os demais números?
Ajustá-los contra o conjunto anotado é prática comum, e é também o que impede levar
um método a outro corpus. Adotamos um critério simples --- \emph{o parâmetro
precisa das respostas certas para ser fixado?} --- e trocamos cada constante por um
estimador medido, durante a construção, sobre o corpus \emph{não} anotado. O corte de rótulo comum vira o quantil $0{,}99$ dos graus
daquele tipo, $h_\kappa=\lceil Q_{0,99}\{\deg(v):v\in\mathcal{S}_\kappa\}\rceil$:
uma contagem absoluta quebraria com a escala, pois num corpus dez vezes maior todo
slot a ultrapassaria e o mecanismo se desligaria sozinho. Pelo mesmo raciocínio, o
limiar de paráfrase vira um quantil alto dos cossenos realmente observados --- é
propriedade do codificador, não da tarefa --- e o filtro de ator vira
$\pi=1/|\Sigma|$ e $\Omega=\mathrm{mediana}_\varepsilon\max_a\omega(a,\varepsilon)$,
apertando-se sozinho quando há mais falantes. Já as palavras que só servem de
moldura (``conversa'', ``data'', ``tipo'') passam a ser detectadas por um contraste
de frequências, o análogo do IDF do lado da pergunta: uma palavra de tópico é
frequente nas perguntas \emph{porque} o é nas conversas, ao passo que uma de
moldura é frequente nas perguntas e ausente do que as pessoas disseram. Os pesos $w_\kappa$, $\lambda$ e $\gamma$ seguem fixos de propósito: afirmam algo
sobre o \emph{modelo} --- que um casamento de conceito informa mais do que um
palpite de hiperônimo ---, não sobre o corpus.
 no corpo do artigo.
% Pacotes assumidos: amsmath, amssymb.
% =====================================================================

\section{MEST: Memória Episódica de Slots Tipados}
\label{sec:metodo}

Chamamos de \textbf{MEST} (\emph{Memória Episódica de Slots Tipados}) o método
desta seção: dado um histórico de conversas em sessões e uma pergunta em
linguagem natural, ele seleciona turnos \emph{verbatim} que caibam em um
orçamento fixo de \emph{tokens} e os entrega ao gerador. Duas restrições governam
o desenho --- a construção da memória não faz nenhuma chamada a modelo de
linguagem, e nenhum parâmetro é fixado a partir de rótulos de ouro.

\subsection{Diagnóstico e princípio de projeto}
\label{sec:metodo:diagnostico}

Partimos de uma medição sobre a linha de base de grafo de incidência
entidade--turno. Apenas $0{,}5\%$ dos turnos de evidência estão ausentes do grafo
--- a cobertura é total ---, mas, entre os turnos de evidência \emph{não
recuperados}, a fração que compartilha algum vértice de entidade com a pergunta é
de $13\%$ (\emph{single-hop}), $7\%$ (\emph{multi-hop}), $10\%$ (temporal) e $5\%$
(\emph{open-domain}). De $87\%$ a $95\%$ dos erros são turnos para os quais o
grafo não oferece caminho algum: não é falha de ranqueamento, porque não há o que
ranquear. Quatro pontes ausentes se repetem nesses casos, e cada uma vira uma
espécie de vértice. \emph{A resposta}: a evidência costuma ser o turno que carrega
o valor e nenhum termo do tópico, que ficou no turno acima --- logo a unidade
atômica do diálogo é o par adjacente, não o turno. \emph{O predicado}: a pergunta
é enquadrada por um verbo (``o que James \emph{adotou}''), invisível a um índice
de substantivos. \emph{O tipo}: pergunta-se por categoria (``quais
\emph{comidas}'') e responde-se por instância (``\emph{frango assado}'').
\emph{A pessoa}: de $98{,}5\%$ a $99{,}7\%$ das perguntas nomeiam exatamente um
falante, e a evidência é dele em $96\%$ a $100\%$ dos casos. Daí o princípio do
método: uma pergunta de memória conversacional é uma \emph{tupla com um buraco}
--- (quem, fez-o-quê, com-o-quê, quando) com um campo vazio --- e a memória deve
ser indexada no mesmo vocabulário em que a pergunta é formulada.

\subsection{Substrato e ingestão}
\label{sec:metodo:substrato}

Cada conversa vira um grafo com rotação (\emph{fatgraph})
$G=(V,H,\alpha,\sigma)$, com $H$ as meias-arestas, $\alpha$ o emparelhamento que
as agrupa em arestas e $\sigma$ a permutação que fixa, em cada vértice, uma ordem
cíclica das incidências. O conjunto de vértices é bipartido,
$V=\mathcal{E}\uplus\mathcal{S}$: em $\mathcal{E}$, os \emph{episódios} --- turnos
contíguos e topicamente coesos ---; em $\mathcal{S}=\biguplus_{\kappa\in K}
\mathcal{S}_\kappa$, os \emph{slots}, tipados por $K=\{\textsf{ator},
\textsf{predicado},\textsf{conceito},\textsf{tipo},\textsf{tempo}\}$, isto é, as
quatro pontes acima mais o conceito, que a linha de base já indexava. Uma aresta
$\varepsilon\!-\!v$ é uma \emph{incidência observada} --- $v$ foi literalmente
pronunciado nos turnos de $\varepsilon$ ---, de modo que $\deg(v)$ é frequência
real e não artefato de parafraseamento.

A rotação é projetada, não herdada da ordem de processamento. Em um vértice de
slot, $\sigma$ é cronológica: a órbita $\mathrm{Orb}(v)$ é ``tudo o que esta
pessoa, este verbo ou este tópico tocou, em ordem'', obtida por leitura e não por
ranqueamento. Em um vértice de episódio, $\sigma$ segue a ordem fixa
$\textsf{ator}\prec\textsf{pred}\prec\textsf{conc}\prec\textsf{tipo}\prec
\textsf{tempo}$, e assim seus \emph{cantos} --- pares de incidências consecutivas
--- são (quem, fez-o-quê), (fez-o-quê, com-o-quê) e (com-o-quê, quando): uma
pergunta é um canto com um lado em branco. As \emph{faces} não são usadas, pois
uma só delas absorve milhares de meias-arestas. A ingestão é determinística: uma
passagem de análise sintática por turno produz sintagmas nominais, lemas verbais,
menções a pessoas e datas de uma vez. Um episódio começa em fronteira de sessão,
cola incondicionalmente os $m_{\min}$ primeiros turnos --- é o par adjacente, e
uma resposta \emph{não} sobrepõe vocabulário com a pergunta justamente quando a
responde --- e admite os seguintes enquanto a coesão de conceitos se mantiver.
Conceitos são elevados a hiperônimos WordNet; datas relativas são resolvidas
contra a data da sessão e indexadas em \emph{todos} os níveis que suportam (ano,
mês, dia).

\subsection{Consulta: a pergunta no vocabulário da memória}
\label{sec:metodo:consulta}

A pergunta $q$ é analisada pelo mesmo extrator, produzindo o conjunto
$\Lambda(q)\subseteq\mathcal{S}$ dos slots que resolvem em vértices existentes.
Cada slot ligado contribui com um peso amortecido pelo grau, e o escore de um
episódio soma canais estruturais e canal denso:
\begin{equation}
\begin{aligned}
\psi(v)&=w_{\kappa(v)}\bigl(1+\ln \deg(v)\bigr)^{-\gamma}\ \ \text{se}\ \deg(v)<h_{\kappa(v)};\quad
\psi(v)=w_{H}\ \text{caso contrário (\emph{hub})},\\
s(\varepsilon)&=w_{D}\,\delta(q,\varepsilon)+\sum_{v\in\Lambda^{*}(q,\varepsilon)}\psi(v)
+\beta\,\mathbb{1}\bigl[\varepsilon\in\mathrm{Orb}(v^{\star})\bigr],
\end{aligned}
\label{eq:psi}
\end{equation}
em que $\Lambda^{*}(q,\varepsilon)$ são os slots não-\textsf{ator} de $\Lambda(q)$ cuja
órbita contém $\varepsilon$, e $\delta(q,\varepsilon)$ é o maior cosseno entre a pergunta
e o episódio ou qualquer um de seus turnos. A intuição do amortecimento é direta: um slot incidente a $200$
episódios não distingue nenhum deles, enquanto um incidente a três carrega a
consulta inteira; acima do corte $h_\kappa$ ele deixa de ser enumerado e vira
bônus constante --- ``\emph{um hub é filtro, nunca ponte}''. O corte é \emph{por
espécie}, pois as escalas de grau são incomparáveis por construção: um ator incide
em metade dos episódios e um conceito, em poucos. É também o amortecimento que
dispensa qualquer parâmetro de granularidade temporal, já que o vértice de ano é
praticamente apagado e o de dia recebe peso quase pleno. O termo $\beta$ é a
\textbf{enumeração por órbita}, ativada quando a pergunta pede uma lista (decidido
pela redação, nunca pela categoria anotada): $v^{\star}$ é o slot específico
\emph{mais raro} entre os ligados e toda a sua órbita entra, pois a resposta é uma
lista e a rotação já \emph{é} a lista.

O ator, por sua vez, é uma \emph{partição}: somar uma constante por ator nomeado
deslocaria metade do grafo pelo mesmo valor e não distinguiria nada. A medição da
Seção~\ref{sec:metodo:diagnostico} diz quais episódios são \emph{elegíveis}, e por
isso se aplica multiplicativamente, $\hat{s}(\varepsilon)=s(\varepsilon)\,
[\pi+(1-\pi)\min(\omega(a,\varepsilon)/\Omega,1)]$, com $\omega(a,\varepsilon)$ a
fração do conteúdo do episódio contribuída por $a$, $\pi$ o piso do prior e
$\Omega$ o ponto em que ``este episódio é dele'' já é verdade --- como todo
episódio contém um turno de cada falante, ``$a$ falou aqui'' não discrimina, mas
``$a$ disse a maior parte'' discrimina. A estrutura de cantos fornece ainda um
sinal determinístico de \textbf{abstenção}: com $A_a=\{\varepsilon:
\omega(a,\varepsilon)\ge\tau\}$ e $\Lambda_{\mathrm{esp}}(q)$ os slots específicos
ligados e não-\emph{hub}, $\mathrm{supp}(q)=\max_{v}|\mathrm{Orb}(v)\cap A_a|/
|\mathrm{Orb}(v)|$ distingue \emph{slot ausente} (nada do conteúdo existe na
memória) de \emph{canto vazio} (o conteúdo existe, mas nunca em episódio dessa
pessoa) --- é o que uma pergunta adversarial produz, por nomear uma combinação que
nunca ocorreu. O sinal é sempre reportado, mas mantido \emph{inativo} nesta
configuração: é o único mecanismo capaz de apagar uma resposta correta, e sua taxa
de falsos positivos é medida antes de ser ligada.

\subsection{Emissão e calibração derivada do corpus}
\label{sec:metodo:emissao}

A distinção entre \emph{unidade de índice} e \emph{unidade de emissão} é onde a
aposta do método se realiza. Ao mesmo orçamento, emitir episódios inteiros põe
$18{,}9$ unidades no contexto contra $58{,}7$ da linha de base e perde tudo o que
precisa de amplitude; emitir um turno por episódio em rodadas recupera a amplitude
e derruba o \emph{recall} de \emph{single-hop} de $0{,}858$ para $0{,}682$, porque
o turno que \emph{responde} é quase sempre o parceiro do turno que \emph{casa}. A
forma final pontua turnos, deixando o episódio ser apenas o que os torna
encontráveis:
\begin{equation}
s(t)=\hat{s}\bigl(\varepsilon(t)\bigr)+w_{D}\cos(q,x_t)
+\lambda\,w_{D}\!\!\max_{t'\in\varepsilon(t),\,t'\neq t}\!\!\cos(q,x_{t'}),
\label{eq:turno}
\end{equation}
cujo último termo é a regra da resposta escrita como aritmética: o turno herda uma
fração $\lambda$ da melhor similaridade entre seus irmãos, e o par sobe junto sem
que nada precise ser agrupado. Tomam-se os turnos em ordem decrescente até esgotar
o orçamento.

Resta de onde vêm os limiares. Varrê-los contra o conjunto anotado é prática comum
e é também o que impede transferir o método a outro corpus; adotamos o critério
\emph{``o parâmetro precisa dos rótulos de ouro para ser fixado?''} e substituímos
cada literal por um estimador medido, em tempo de construção, sobre o grafo não
anotado. O corte de \emph{hub} vira $h_\kappa=\lceil Q_{0{,}99}\{\deg(v):v\in
\mathcal{S}_\kappa\}\rceil$: uma contagem absoluta é errada sob mudança de escala,
pois num corpus dez vezes maior todo slot a cruzaria e o mecanismo se desligaria
sozinho (aqui, o corte derivado para conceitos é $16$ contra $60$ absoluto). O
limiar de paráfrase entre conceitos vira um quantil alto dos cossenos observados,
por ser propriedade do codificador tanto quanto da tarefa, e o prior de ator vira
$\pi=1/|\Sigma|$ e $\Omega=\mathrm{mediana}_\varepsilon\max_a\omega(a,\varepsilon)$,
de sorte que nomear um entre oito falantes o fortalece sozinho. A lista de
palavras de enquadramento (``conversa'', ``data'', ``tipo'') deixa de ser manual e
vira $\{w:\mathrm{df}_q(w)\ge d_0 \wedge \mathrm{df}_q(w)/\max(\mathrm{df}_m(w),
m_0)\ge\rho\}$, o análogo do IDF do lado da pergunta: uma palavra de tópico é
frequente nas perguntas \emph{porque} é frequente nas conversas, e sua razão fica
próxima de $1$; uma palavra de \emph{template} é frequente nas perguntas e ausente
da memória. Já os pesos $w_\kappa$, a fração $\lambda$ e o expoente $\gamma$ não
são derivados de propósito: codificam uma afirmação sobre o \emph{modelo} --- ``um
casamento de conceito informa mais que um palpite de hiperônimo'' --- e não sobre
o corpus. Declaramos, por fim, que o estimador de enquadramento é
\emph{transdutivo}, pois lê o texto das perguntas, embora nunca as respostas, a
evidência ou a categoria.
